% ============================================================================
\section{Latar belakang dan rumusan masalah}
\label{sec:masalah}
% ============================================================================

Bagian ini memberi pembaca dari luar Indonesia pengetahuan domain secukupnya untuk memahami mengapa perubahan PPh 21 sulit diikuti, lalu merumuskan masalahnya.

\subsection{Perhitungan PPh 21 pegawai tetap}
\label{subsec:pph21}

PPh 21 pegawai tetap dihitung dari penghasilan bruto. Bruto terdiri atas penghasilan teratur, seperti gaji, tunjangan, dan uang lembur, dan penghasilan tidak teratur, seperti tunjangan hari raya (THR) dan bonus. Premi jaminan kecelakaan kerja (JKK), jaminan kematian (JKM), dan jaminan kesehatan yang dibayar pemberi kerja ikut menjadi bruto, sedangkan iuran jaminan hari tua (JHT) dan jaminan pensiun (JP) yang dibayar pemberi kerja tidak~\cite{pmk168_2023}. Penghasilan neto adalah bruto dikurangi biaya jabatan, yaitu 5\% dari bruto dengan batas Rp500.000 sebulan dan Rp6.000.000 setahun, serta iuran JHT dan JP yang dibayar pegawai. Iuran JP pegawai sebesar 1\% dihitung dari upah paling tinggi sebesar batas upah yang ditetapkan setiap tahun. Penghasilan kena pajak (PKP) adalah neto setahun dikurangi penghasilan tidak kena pajak (PTKP), yaitu Rp54 juta untuk diri wajib pajak, ditambah Rp4,5 juta bila kawin dan Rp4,5 juta untuk setiap tanggungan sampai tiga orang~\cite{pmk101_2016}. PKP dibulatkan ke bawah hingga ribuan rupiah. PPh setahun dihitung dengan tarif progresif Pasal 17 UU PPh, yang sejak UU HPP terdiri atas lima lapisan, yaitu 5\% sampai Rp60 juta, 15\% sampai Rp250 juta, 25\% sampai Rp500 juta, 30\% sampai Rp5 miliar, dan 35\% di atasnya~\cite{uu7_2021}.

Cara memotong PPh 21 setiap bulan berubah pada 2024. Sampai tahun pajak 2023, pemberi kerja menyetahunkan penghasilan teratur setiap bulan, menghitung PPh setahun, lalu membaginya per bulan. Penghasilan tidak teratur dipajaki terpisah sebagai selisih PPh setahun dengan dan tanpa penghasilan itu~\cite{per16_2016}. Sejak 2024, potongan bulan Januari sampai November adalah TER dikali bruto bulan itu. TER terbagi menjadi kategori A, B, dan C menurut status PTKP, dengan 44, 40, dan 41 lapisan tarif bulanan~\cite{pp58_2023}. Pada masa pajak terakhir, yaitu Desember atau bulan pegawai berhenti bekerja, pemberi kerja menghitung PPh setahun dengan tarif Pasal 17, lalu memotong selisihnya dengan jumlah yang sudah dipotong~\cite{pmk168_2023}. Selisih ini dapat negatif. Dalam hal itu, kelebihan potongan dikembalikan kepada pegawai.

Pemberi kerja dapat memotong PPh 21 dari penghasilan pegawai, memberi tunjangan pajak sebesar PPh 21 itu sendiri (\en{gross-up}), atau menanggungnya. Pada \en{gross-up}, tunjangan pajak adalah PPh atas bruto yang memuat tunjangan itu, sehingga nilainya didefinisikan secara melingkar. Pada 2025 dan 2026, pemerintah juga menanggung PPh 21 pegawai di sektor tertentu yang bruto tetap dan teraturnya paling banyak Rp10 juta sebulan. PPh yang ditanggung pemerintah dibayarkan tunai kepada pegawai dan tidak dihitung sebagai penghasilan~\cite{pmk10_2025,pmk72_2025,pmk105_2025}. Tabel~\ref{tab:ruang-pajak} merangkum dimensi yang menentukan besarnya pajak.

\begin{table}[t]
\centering
\caption{Dimensi perhitungan PPh 21 pegawai tetap yang dicakup}
\label{tab:ruang-pajak}
\scriptsize
\begin{tabularx}{\columnwidth}{@{}P{2.0cm} L@{}}
\toprule
\textbf{Dimensi} & \textbf{Nilai} \\
\midrule
Rezim & Disetahunkan (sampai 2023); TER (sejak 2024); TER dengan fasilitas DTP (2025--2026) \\
Lapisan Pasal 17 & 5\%, 15\%, 25\%, 30\%, 35\% (empat lapisan sebelum UU HPP, hanya untuk data pengembangan) \\
Kategori TER & A, B, C dengan 44, 40, dan 41 lapisan tarif bulanan \\
Status PTKP & TK/0--TK/3, K/0--K/3; menentukan kategori TER \\
Metode & Dipotong (\en{gross}), \en{gross-up}, ditanggung pemberi kerja \\
Pola kerja & Setahun penuh, masuk di tengah tahun, keluar di tengah tahun \\
Penghasilan & Teratur; tidak teratur (THR, bonus) \\
\bottomrule
\end{tabularx}
\end{table}

\status{Jumlah lapisan TER 44, 40, dan 41 berasal dari \texttt{dataset/01\_regulasi/tables/ter\_bulanan.csv}, yang sudah diekstraksi dua kali, dicek visual pada beberapa halaman lampiran PP 58/2023, dan cocok dengan contoh PMK 168/2023. Catatan revisi masih meminta cek ulang penuh ke lampiran PP 58/2023 (Tahap 2).}

Penelitian ini terbatas pada pegawai tetap yang menerima penghasilan dari satu pemberi kerja, pada tahun pajak 2023--2026. Pegawai tidak tetap, bukan pegawai, pensiunan, pegawai dengan lebih dari satu pemberi kerja, pindah cabang, fasilitas khusus Ibu Kota Nusantara, dan penilaian natura berada di luar lingkup. Nilai natura diterima sebagai masukan. Dimensi di Tabel~\ref{tab:ruang-pajak} saling membatasi. Status PTKP menentukan kategori TER. Metode disetahunkan hanya berlaku sampai 2023, TER hanya sejak 2024, dan DTP hanya pada 2025--2026 untuk sektor dan batas bruto tertentu. Pegawai yang masuk atau keluar di tengah tahun mengubah masa pajak terakhir dan cara menghitung PPh setahun. Ruang uji di Bagian~\ref{subsec:ruanguji} hanya membangkitkan kombinasi yang sah menurut batasan ini.

\terbuka{Prototipe saat ini menghitung \en{gross-up} hanya pada rezim TER. Apakah \en{gross-up} pada rezim disetahunkan (2023) masuk lingkup, atau dicatat sebagai batasan antardimensi?}

\subsection{Kebijakan perusahaan dan batas kewenangannya}
\label{subsec:kebijakan}

Perusahaan menetapkan sebagian besar komponen penghasilan pegawai lewat peraturan perusahaan, perjanjian kerja bersama (PKB), atau keputusan direksi. Contohnya cara memprorata gaji pegawai yang masuk di tengah bulan, cara menghitung THR, besaran tunjangan, dan dasar upah iuran BPJS. Kebijakan seperti ini menentukan besarnya bruto, sehingga ikut menentukan PPh 21. Perusahaan juga boleh memilih di antara pilihan yang disediakan regulasi, misalnya menanggung PPh 21 pegawai atau memberinya tunjangan pajak.

Perusahaan tidak boleh menyimpangi ketentuan wajib. Contohnya klasifikasi jenis penghasilan. PER-16/PJ/2016 Pasal 1 angka 15 dan PMK 168/2023 Pasal 5 ayat (3) huruf a menggolongkan uang lembur sebagai penghasilan teratur~\cite{per16_2016,pmk168_2023}. Kebijakan payroll Perusahaan X mencatatnya sebagai penghasilan tidak teratur. Pada metode disetahunkan, penghasilan teratur ikut disetahunkan sedangkan penghasilan tidak teratur tidak, sehingga perbedaan klasifikasi ini mengubah besarnya potongan bulanan. Bila kebijakan perusahaan bertentangan dengan ketentuan wajib, ketentuan regulasi yang berlaku (\en{lex superior}), dan konflik itu harus tercatat. Di \en{spreadsheet} Perusahaan X, konflik ini tersembunyi di dalam rumus.

\subsection{Jenis perubahan regulasi}
\label{subsec:jenis-perubahan}

Tabel~\ref{tab:jenis} memuat taksonomi jenis perubahan yang dipakai untuk menyusun korpus dan melaporkan hasil. Satu peristiwa dapat termasuk lebih dari satu jenis. Contohnya PMK 10/2025, yang menambah fasilitas baru sekaligus berlaku surut. Saat melaporkan tingkat keberhasilan, jenis yang mudah seperti perubahan parameter tidak digabung dengan jenis struktural seperti pergantian metode.

\begin{table}[t]
\centering
\caption{Taksonomi jenis perubahan}
\label{tab:jenis}
\scriptsize
\begin{tabularx}{\columnwidth}{@{}P{1.8cm} P{2.5cm} L@{}}
\toprule
\textbf{Jenis} & \textbf{Contoh} & \textbf{Wujud pada \en{knowledge base}} \\
\midrule
Parameter & Batas upah JP setiap Maret & Nilai parameter baru dengan masa berlaku baru \\
Tabel tarif & Tabel TER (PP 58/2023) & Versi tabel baru \\
Klasifikasi & Perlakuan natura dan kenikmatan sebagai penghasilan (PMK 66/2023) & Pemetaan jenis penghasilan ke kategori pajak berubah \\
Aturan atau fasilitas baru & DTP padat karya (PMK 10/2025) & Aturan baru, sering disertai isian data baru \\
Kondisi atau cakupan & Sektor pariwisata (PMK 72/2025) & Kondisi aturan yang ada diperluas atau dipersempit \\
Pergantian metode & Disetahunkan ke TER (2024) & Sekelompok aturan diganti, aturan lama diakhiri \\
Pencabutan & PMK 250/2008 dicabut PMK 168/2023 & Masa berlaku aturan diakhiri secara eksplisit \\
Berlaku surut & PMK 10/2025 (terbit Februari, berlaku Januari) & Tanggal mulai berlaku lebih awal dari tanggal terbit \\
Kebijakan perusahaan & Versi PKB baru & Aturan lapisan perusahaan berubah \\
\bottomrule
\end{tabularx}
\end{table}

\status{Taksonomi ini adalah usulan draf yang menambah satu jenis dari daftar minimum di \en{skeleton}, yaitu perubahan kondisi atau cakupan. Taksonomi ditetapkan final di Tahap 2, bersama taksonomi konflik dan taksonomi kesalahan LLM.}

\subsection{Rumusan masalah}
\label{subsec:rumusan}

Masukan masalah ini adalah \en{knowledge base} $\K_{n-1}$ yang benar untuk keadaan hukum sebelum peristiwa perubahan ke-$n$, beserta dokumen sumber peristiwa itu. Keluarannya adalah $\K_n$ yang benar untuk keadaan hukum sesudahnya, dengan tiga syarat. Kode \en{engine} tidak berubah. Hasil yang tidak seharusnya berubah tetap sama, termasuk seluruh hasil untuk masa sebelum perubahan berlaku. Kebijakan perusahaan yang menjadi bertentangan dengan ketentuan wajib dilaporkan. Masalah ini berulang, karena $\K_n$ menjadi masukan untuk peristiwa berikutnya. Kesalahan yang lolos pada satu peristiwa dapat terbawa, atau memicu kesalahan baru, pada peristiwa sesudahnya. Bagian~\ref{subsec:adaptif} menyatakan syarat-syarat ini secara formal.
